\section{The Tessellation Survey}

The geometry of the mesh is fixed by selection, not by taste. To make that selection honest one cannot examine a single candidate and declare it good. One must run the whole enumerated family of regular hyperbolic tessellations through the identical battery of measurements, lay the numbers side by side, and see which substrate the comparison forces. This section reports that cross-substrate survey. The other tessellations appear here precisely as the comparison that makes the choice of $\{3,4,3,4\}$ forced rather than assumed.

The survey rests on a fact that makes the search tractable. The regular hyperbolic tessellations are a finite, fully enumerated set. Compact regulars stop at four dimensions. Regulars of any class stop at five. There is nothing beyond. So the catalog is a closed list, every member can be measured, and the comparison is complete rather than a sample of an open-ended space.

\begin{principle}[The selection is over a finite catalog]
The regular hyperbolic tessellations form a bounded, fully enumerated set. The compact regulars terminate in four dimensions and the regulars of every class terminate in five. The survey can therefore exhaust the catalog, and the selection of the mesh is a choice among finitely many fixed candidates.
\end{principle}

\subsection{One battery, one module, every tessellation}

The strength of the survey is uniformity. Every substrate is measured the same way, so the numbers are directly comparable across dimensions and classes. A single reusable module takes a tessellation's Schl\"afli symbol and measures it. It builds the Coxeter reflection-group mesh, which is uniform across all dimensions and classes, so a two-dimensional tiling and a five-dimensional pentacomb are run through the identical procedure.

For each substrate the module records four properties.

\begin{itemize}
\item \textbf{Curvature.} The shell growth ratio (how fast the count of docks grows shell by shell) together with the Gram signature of the substrate. These read off the sign and strength of the curvature from the inside.
\item \textbf{Crystallographic character.} Whether the substrate can carry a root system. This is the gauge hook, the structural opening through which an emergent gauge symmetry can appear.
\item \textbf{Spinor sites.} Whether the $24$ sites out of a dock carry the $24$-cell and $D_4$ spinor sectors, the home of the $8s$ and $8c$ pieces. This is the spin hook.
\item \textbf{Propagation.} Whether a clean K\"ahler-Dirac fermion stays extended and moves across the mesh, or localizes and stays put.
\end{itemize}

A short reminder of the vocabulary is worth giving before the results. The \emph{mesh} is the whole tessellation, a weave of docks. A \emph{dock} is one geometric cell, a here-now. A \emph{site} is one of the $24$ directions out of a dock, the directed slot in which a tone lives. The \emph{tone} is the ternary value $\{-1, 0, +1\}$ a site carries, read as pain, peace, pleasure. The law that advances the state each beat is the \emph{knit}, collide then stream. None of these change across the survey. What changes is the shape of the mesh, the Schl\"afli symbol, and the survey asks which shape carries the spinor sites natively.

\subsection{The catalog}

The catalog holds $45$ regular tessellations. Of these, $42$ are buildable by the integer Coxeter engine. The three that are not buildable are cataloged only for completeness. They are the star regulars with fractional Schl\"afli symbols, such as $\{7/2, 7\}$ and $\{3, 3, 5, 5/2\}$, and the apeirogonal regulars at infinity, such as $\{\infty, 3\}$. The $42$ buildable substrates span four dimensions.

\begin{table}[ht]
\centering
\begin{tabular}{@{}lcccc@{}}
\toprule
dimension & compact & paracompact & noncompact & buildable total \\
\midrule
2D & 10 (sampled family) & 0 & 0 & 10 \\
3D & 4 (incl. $\{5,3,4\}$) & 11 & 3 & 18 \\
4D & 5 & 2 (the 24-cell pair) & 2 & 9 \\
5D & 0 & 5 (pentacombs) & 0 & 5 \\
\bottomrule
\end{tabular}
\caption{The buildable regular hyperbolic tessellations by dimension and class.}
\label{tab:tessellation-catalog}
\end{table}

The two-dimensional row is a named sample of an infinite family of plane tilings. It is the testbed plane, the cheapest place to check a measurement. Three dimensions is where the nearest neighbor $\{5,3,4\}$ lives, the twelve-direction substrate that an earlier round of the theory committed to. Four dimensions is where the $24$-cell first appears. Five dimensions is the last dimension with any regular honeycomb at all, the five pentacombs, and there the catalog ends.

\subsection{Result one, matter is universal}

The first result is that matter is not picky about the substrate.

\begin{result}[Matter is universal]
A clean K\"ahler-Dirac fermion propagates on every one of the $42$ buildable regular hyperbolic tessellations. The propagating, extended phase is universal across the whole catalog.
\end{result}

Every one of the $42$ is hyperbolic. The Gram signature is Lorentzian in each case, the growth ratio exceeds one in each case, and the growth ratio rises with dimension. The clean fermion sits in the extended, propagating phase on every substrate tested, with the clean return probability in a tight band from $0.048$ to $0.063$ across the entire catalog. The control is a strong deterministic disorder that localizes the fermion, driving the localized return probability up from about $0.31$ past $0.95$. There is no randomness anywhere in this. The base is deterministic, and the disorder is a fixed deterministic pattern, not a random one.

\begin{table}[ht]
\centering
\begin{tabular}{@{}lc@{}}
\toprule
dimension & growth ratio (range) \\
\midrule
2D & about $1.17$ to $1.50$ \\
3D & about $1.32$ to $1.63$ \\
4D & about $1.38$ to $1.95$ \\
5D & about $1.64$ to $1.83$ \\
\bottomrule
\end{tabular}
\caption{The shell growth ratio rises with dimension across the catalog. Every value exceeds one, the signature of negative curvature.}
\label{tab:growth-ratio}
\end{table}

The reason a hyperbolic geometry can host a propagating fermion is the negative curvature itself, which every one of the $42$ shares. No special direction structure is needed for matter to move. The exponential room of a hyperbolic mesh is what carries the fermion outward, and that room is present on every member of the catalog.

\emph{Matter is not special to $\{5,3,4\}$ or to $\{3,4,3,4\}$. The fermion lives and moves on any hyperbolic substrate.}

\subsection{Result two, the spinor sites are rare and dimension-gated}

The second result is the sharp one, and it is the opposite of the first. Where matter is universal, the spinor sites are rare.

\begin{result}[The spinor sites are rare]
The native $24$-cell and $D_4$ spinor sites appear in exactly seven of the $42$ buildable tessellations. They are the $24$-cell-faceted ones, two in four dimensions and five in five dimensions, and none below.
\end{result}

The $24$-cell and $D_4$ sites, the $[3,4,3]$ subdiagram that houses the $8s$ and $8c$ spinor sectors, appear only in the $24$-cell-faceted tessellations. The seven, listed by dimension, are these.

\begin{table}[ht]
\centering
\begin{tabular}{@{}ll@{}}
\toprule
dimension & spinor-site substrates \\
\midrule
2D & none \\
3D & none \\
4D & $\{3,4,3,4\}$, $\{4,3,4,3\}$ (the 24-cell pair) \\
5D & $\{4,3,3,4,3\}$, $\{3,4,3,3,4\}$, $\{3,3,4,3,3\}$, $\{3,4,3,3,3\}$, $\{3,3,3,4,3\}$ \\
\bottomrule
\end{tabular}
\caption{The seven substrates carrying native spinor sites. Two in four dimensions, five in five dimensions, none below four.}
\label{tab:spinor-sites}
\end{table}

Two in four dimensions plus five in five dimensions gives seven in total, with none in two or three dimensions. The feature is dimension-gated. It cannot appear until four dimensions, the dimension where the $24$-cell first exists, and it vanishes again past five, where the catalog itself ends.

These same seven are exactly the substrates that are crystallographic and spinor-carrying together. The $24$-cell sites are three-fold and four-fold, so the spinor sites bring the gauge hook with them. One feature delivers both gifts at once. There is no member of the catalog that has the spinor sites without also being crystallographic, and the survey makes that coincidence explicit rather than assuming it.

\subsection{The picture this gives}

Laying the two results side by side gives the shape of the whole survey. One property is universal, one is rare, and the survey separates them cleanly.

\begin{table}[ht]
\centering
\begin{tabular}{@{}lll@{}}
\toprule
property & how common & which \\
\midrule
hyperbolic, hosts a propagating fermion & universal (42 of 42) & every regular hyperbolic tessellation \\
crystallographic (gauge hook) & common (21 of 42) & the 3-, 4-, 6-fold ones \\
the 24-cell / $D_4$ spinor sites & rare (7 of 42) & $\{3,4,3,4\}$, $\{4,3,4,3\}$, the five pentacombs \\
spinor sites and curvature and gauge together & rare (7 of 42) & the same seven \\
\bottomrule
\end{tabular}
\caption{The universal-versus-rare separation. Matter is everywhere, the spinor sites are the dimension-gated rarity.}
\label{tab:rare-vs-universal}
\end{table}

The survey separates a universal property from a rare one. Matter, a propagating fermion, is everywhere in hyperbolic geometry. The native spinor sites, with their gauge hook riding along, are the rare and dimension-gated feature. That rare feature is exactly what the $\{3,4,3,4\}$ commitment and the pentacomb resolution are about. The detail of why four dimensions and the sites select $\{3,4,3,4\}$ is taken up where the dimension ladder is climbed.

\begin{principle}[Matter is common, the spinor sites are the discriminator]
Hosting a propagating fermion is universal across the catalog and so cannot select any substrate. The native $24$-cell spinor sites appear in only seven substrates, and they are what a selection argument can grip. The geometry decision is therefore a decision about the spinor sites, not about whether matter can move.
\end{principle}

\subsection{Where $\{5,3,4\}$ sits}

The earlier three-dimensional pick, $\{5,3,4\}$, sits exactly where the survey predicts. It is hyperbolic. It hosts a propagating fermion like all the others, with a clean return probability of $0.0508$, squarely in the universal-matter class. But it is not in the rare-sites class. Its twelve directions do not carry the spinor sectors. Under the icosahedral symmetry those twelve directions decompose with no spinor anywhere in the list, whereas the $24$ sites of a $24$-cell-faceted substrate decompose as $8v + 8s + 8c$ and so carry spinors directly.

So $\{5,3,4\}$ must get its spin from the form complex or from a defect, not from the sites themselves. That is the whole story of why $\{5,3,4\}$ needs extra machinery for spin, and why a pentacomb, a sites-carrying curved substrate, resolves the trade. A substrate that carries the spinor sites natively pays no such cost.

\subsection{Honest scope}

The survey is a comparative measurement on the Coxeter reflection-group mesh, the uniform model that lets every tessellation be measured the same way. Its claims are scoped accordingly.

\begin{itemize}
\item The growth ratio is the \emph{measured} curvature signature, read off the shell counts.
\item The spinor and crystallographic hooks are \emph{exact} properties of the Schl\"afli symbol, not measured quantities but combinatorial facts.
\item The propagation is the extended-versus-localized return-probability test, with the deterministic-disorder control. It is the same test that established the $\{5,3,4\}$ and pentacomb results elsewhere in the paper.
\end{itemize}

This is a catalog comparison. The value it delivers is the universal-versus-rare separation, not a new emergent claim about any one substrate. It does not derive matter from the substrate, it shows that matter is common while the spinor sites are rare, and it locates the few substrates where the rare feature lives. The full per-tessellation data, including the propagation column for all $42$, lives in the companion appendix table.

\subsection{The verified claims}

Three claims come out of the survey verified against the full catalog.

\begin{table}[ht]
\centering
\begin{tabular}{@{}l@{}}
\toprule
claim \\
\midrule
every buildable regular hyperbolic tessellation is hyperbolic and hosts a propagating fermion \\
the $24$-cell / $D_4$ spinor sites appear in exactly the seven $24$-cell-faceted tessellations \\
the spinor-site set equals the crystallographic-and-spinor set \\
\bottomrule
\end{tabular}
\caption{The verified results of the tessellation survey.}
\label{tab:verified-claims}
\end{table}

The survey thus does its one job. It runs the whole enumerated family through one battery and shows that matter is universal while the native spinor sites are rare and arrive only in four dimensions and above. That rarity is the lever the geometry decision pulls, and the next section climbs the dimension ladder to see how the lever lands on $\{3,4,3,4\}$.
